% ==============================================================================
% PAGE 47: CHAPTER-11 CONCLUSION (Numbered Page 39)
% ==============================================================================
\begin{center}
  {\fontsize{15}{18}\selectfont \bfseries \textcolor{red!80!black}{\underline{CHAPTER-11}}}\\[2mm]
  {\fontsize{16}{20}\selectfont \bfseries \textcolor{red!80!black}{\underline{CONCLUSION}}}\\[6mm]
\end{center}

\noindent
\begin{spacing}{1.25}
{\fontsize{10.2}{15}\selectfont
The designed parabolic solar water heater successfully utilizes solar energy for water heating and small-scale power generation.\\[3mm]
The use of a parabolic dish with mirror pieces ensures effective concentration of sunlight on the receiver tube.\\[3mm]
The automatic single-axis tracking system increases the efficiency by keeping the dish aligned with the sun throughout the day.\\[3mm]
A 6V DC motor controlled by LDR sensors enables gradual rotation of about 4--5 degrees per movement for accurate tracking.\\[3mm]
The generated heat is stored in water, providing a renewable and eco-friendly source of energy.\\[3mm]
The attached solar panel also generates power to charge the 6V battery used for operating the tracking system.\\[3mm]
This project demonstrates the practical use of solar energy in both heating and automation applications.\\[3mm]
The system design is simple, cost-effective, and suitable for domestic and educational purposes.\\[3mm]
It helps in reducing dependence on conventional energy sources and promotes sustainable technology.\\[3mm]
Overall, the project meets its objectives of energy conservation, automation, and efficient solar energy utilization.
}
\end{spacing}
\clearpage
